\renewcommand\thesection{\Alph{section}}
\renewcommand*{\theHsection}{appedix.\thesection}
\setcounter{section}{0}
\setcounter{figure}{3}
\setcounter{table}{6}
\setcounter{equation}{4}

\iffalse

A. Qualitative Results
1 Gallery 4 tasks
    each task 5 cases 
2 comparison 
    text-to-motion
        MDM MLD T2MGPT Ours          tm2t
        5个例子
    motion-to-text
        tm2t
        ours
   text descriptions
   
B. Additional Experiments
 B.1 VAE Quantization
 B.2 Text-to-motion -KIT
 B.3 No-pretraining only instruction
    固定迭代次数
 B.4 Details on Comparisons
     Implement details about MDM

 % B.1 Evaluation of Other Language Models τ

C. Discussion on differencet of T2M-GPT
   the difference between our and t2m-gpt
D. Inference time
  FPS 
  
E. 可视化、统计 codebook柱状图 利用率
   512 1024
   
F. Evaluation of Hyperparameters
    Generation number  搜索路径
    

G. User-study
H. Details on MotionGPT Models
    G.1 Loss
    G.2 Network Architectures
        xxxxx
    G.3 Implementation Details
        Small Base large, others xxxx

I. Details on Instruction datseset
J. Metric Definitions
    NLP metrics
    Completion metrics

videos
    入场 华尔兹 从左到右
    1 input - output
        two cases four tasks
    2 comparison of text-to-motion
        two
    3 comparison of motion-to-text
        two
    4 comparison of motion completion
        two
    5. More results
        5~10 例子
    
    结束页say goodbye motion 
\fi

This appendix provides qualitative comparison results (\cref{sec:appendix:qualitative}), 
additional experiments (\cref{sec:appendix:exps}) on the components of MotionGPT models, inference time (\cref{sec:inferencetime}), statistics on motion vocabulary (\cref{sec:appendix:visual}),
evaluations on hyperparameters (\cref{sec:appendix:hyperparameters}), user study (\cref{sec:appendix:userstudy}), a protocol for the uniform evaluation (\cref{sec:appendix:data: details}), and more implementation details (\cref{sec:appendix:model:details}) of MotionGPT models .
Please note evaluations on our training scheme (\cref{sec:pretrian}), elaborations on the difference of T2M-GPT (\cref{sec:appendix:diff}),
implementation details of motion completion (\cref{sec:appendix:reimplement}),
and more metric definitions (\cref{appedix:metrics:details}).

\textbf{Video.} We have provided the supplemental video to illustrate our results. In this video, we show 1) comparisons of text-to-motion, 2) comparisons of motion captioning, and 3) more results on motion prediction and other tasks. We suggest watching this video for dynamic motion results.

\textbf{Code} is available in supplements. We provide example code files, which include the process of the training and evaluation of our MotionGPT models, as well as several example results. We can hardly upload our large model files, but all codes, data, and pre-trained models will be fully released.


% \textbf{Video.} We have provided supplemental videos. In these supplemental videos, we show 1) comparisons of text-based motion generation, 2) comparisons of motion captioning, and 3) more samples of unconstrained generation. We suggest watching this video for dynamic motion results.\\

% \textbf{Code} is available on \href{https://github.com/chenfengye/motion-latent-diffusion}{GitHub Page}. We provide the process of the training and evaluation of MLD models, the pre-trained model files, the demo script, and example results.


\section{Qualitative Results}
\label{sec:appendix:qualitative}
% We visualize some qualitative results on the comparison of text-to-motion ($cf.$ \cref{fig:appendix:compt2m}), motion-to-text ($cf.$ \cref{fig:appendix:compm2t}), and our result gallery on multiple tasks ($cf.$ \cref{fig:appendix:more_rs}).

\begin{figure}[h]
    \centering
    \includegraphics[width=\linewidth]{comp_t2m.pdf}
    \caption{Comparison on text-driven motion generation. The provided state-of-the-art methods are under the same training and inference setting on HumanML3D \cite{Guo_2022_CVPR_humanml3d}. The red words and boxes highlight the misaligned motions. The results demonstrate that our motion-language per-training shows promising text understanding for motion generation.}
    \label{fig:appendix:compt2m}
\end{figure}

\begin{figure}[!htb]
    \centering
    \includegraphics[width=0.99\linewidth]{gallery.pdf}
    \caption{Gallery for the results of our unified MotionGPT. More samples are from our best model for text-to-motion synthesis, motion captioning, and textual question-to-answer task. The supervision of MotionGPT relies on our instruction-based motion-language dataset ($cf.$ \cref{sec:appendix:data: details}) based on previous motion datasets~\cite{Guo_2022_CVPR_humanml3d, AMASS_ICCV2019}. We recommend the dynamic visualization in our supplemental video. }
    \label{fig:appendix:more_rs}
\end{figure}

\begin{figure}[h]
    \centering
    \includegraphics[width=\linewidth]{comp_m2t.pdf}
    \vspace{-10pt}
    \caption{Comparison of the state-of-the-art method on motion captioning task. All provided methods are under the same training and inference setting on HumanML3D~\cite{Guo_2022_CVPR_humanml3d}. The results demonstrate that our text descriptions correspond better to the motion and have correct grammar. The orange words indicate the matching results, while the red marks the incorrect grammar.}
    \label{fig:appendix:compm2t}
\end{figure}

\section{Additional Experiments}
\label{sec:appendix:exps}
We conduct several experiments to continue the evaluations of MotionGPT models. We first evaluate the text-to-motion results on KIT dataset (\cref{sec:t2mkit}). Then we evaluate the hyperparameters of motion tokenizer $\mathcal{V}$ (\cref{sec:vqvae}). After that, we study the effectiveness of the training scheme (\cref{sec:pretrian}). We also provide the elaboration on the difference of T2M-GPT (\cref{sec:appendix:diff}),
implementation details of motion completion (\cref{sec:appendix:reimplement}).


\subsection{Text-to-Motion on KIT dataset.} 
\label{sec:t2mkit}
Following the same procedure on HumanML3D\cite{Guo_2022_CVPR_humanml3d} dataset, We train a 220M MotionGPT base model on the KIT\cite{Plappert2016kit} dataset without any pre-training. We evaluate this model under the same settings of \cite{Guo_2022_CVPR_humanml3d}. Most results are borrowed from their own paper of the benchmark in \cite{Guo_2022_CVPR_humanml3d}. \cref{tab:tm:comp:kit} shows that MotionGPT achieves comparable performance compared to the previous state-of-the-arts.

\begin{table}[h]
\vspace{20pt}
\resizebox{\columnwidth}{!}{%
\begin{tabular}{@{}lcccccccc@{}}
\toprule
\multirow{2}{*}{Methods}&\multicolumn{3}{c}{RPrecision$\uparrow$}&\multicolumn{1}{c}{\multirow{2}{*}{FID$\downarrow$}}&\multirow{2}{*}{MMDist$\downarrow$}&\multirow{2}{*}{Diversity$\rightarrow$}&\multirow{2}{*}{MModality$\uparrow$}\\\cmidrule(lr){2-4}
&\multicolumn{1}{c}{Top1}&\multicolumn{1}{c}{Top2}&\multicolumn{1}{c}{Top3}&\multicolumn{1}{c}{}&&&\\\midrule
Real&
$0.424^{\pm.005}$&
$0.649^{\pm.006}$&
$0.779^{\pm.006}$&
$0.031^{\pm.004}$&
$2.788^{\pm.012}$&
$11.08^{\pm.097}$&
\multicolumn{1}{c}{-}
\\\midrule
% TEMOS\cite{petrovich22temos}&
% $0.353^{\pm.006}$&
% $0.561^{\pm.007}$&
% $0.687^{\pm.005}$&
% $3.717^{\pm.051}$&
% $3.417^{\pm.019}$&
% $10.84^{\pm.100}$&
% $0.532^{\pm.034}$\\
TM2T\cite{chuan2022tm2t}&
$0.280^{\pm.005}$&
$0.463^{\pm.006}$&
$0.587^{\pm.005}$&
$3.599^{\pm.153}$&
$4.591^{\pm.026}$&
$9.473^{\pm.117}$&
${3.292}^{\pm.081}$\\
% T2M\cite{Guo_2022_CVPR_t2m}&
% $0.370^{\pm.005}$&
% $0.569^{\pm.007}$&
% $0.693^{\pm.007}$&
% $2.770^{\pm.109}$&
% $3.401^{\pm.008}$&
% $\underline{10.91}^{\pm.119}$&
% $1.482^{\pm.065}$\\
% MotionDiffuse\cite{zhang2022motiondiffuse}&
% $\boldsymbol{0.417}^{\pm.004}$&
% $\boldsymbol{0.621}^{\pm.004}$&
% $\boldsymbol{0.739}^{\pm.004}$&
% $1.954^{\pm.062}$&
% $\boldsymbol{2.958}^{\pm.005}$&
% $\boldsymbol{11.10}^{\pm.143}$&
% $0.730^{\pm.013}$\\
MDM\cite{mdm2022human}&
$0.164^{\pm.004}$&
$0.291^{\pm.004}$&
$0.396^{\pm.004}$&
${0.497}^{\pm.021}$&
$9.191^{\pm.022}$&
${10.85}^{\pm.109}$&
${1.907}^{\pm.214}$\\
%ours&
%$0.383^{\pm.006}$&$0.591^{\pm.007}$&$0.714^{\pm.007}$&$0.730^{\pm.053}$&$3.214^{\pm.034}$&$10.523^{\pm.115}$&$2.394^{\pm.073}$
%\\
MLD\cite{chen2023mld}&
${0.390}^{\pm.008}$&
${0.609}^{\pm.008}$&
${0.734}^{\pm.007}$&
${0.404}^{\pm.027}$&
${3.204}^{\pm.027}$&
$10.80^{\pm.117}$&
${2.192}^{\pm.071}$\\
T2M-GPT \cite{zhang2023generating}    &$0.416^{\pm.006}$ & $0.627^{\pm.006}$ & $0.745^{\pm.006}$ & $0.514^{\pm.029}$ & $3.007^{\pm.023}$ & $10.92^{\pm.108}$ & $1.570^{\pm.039}$  \\
% ReMoDiffuse \cite{remodiffuse} & $\boldsymbol{0.427}^{\pm.014}$ & $\boldsymbol{0.641}^{\pm.004}$ & $\boldsymbol{0.765}^{\pm.055}$ & $\boldsymbol{0.155}^{\pm.006}$ & $\boldsymbol{2.814}^{\pm.012}$ & $10.80^{\pm.105}$  & $1.239^{\pm.028}$\\
\midrule
MotionGPT (Ours) &
$0.366^{\pm.005}$ & $0.558^{\pm.004}$ & $0.680^{\pm.005}$ & $0.510^{\pm.016}$ & $3.527^{\pm.021}$ & $10.35^{\pm.084}$ & 
$2.328^{\pm.117}$
\\\bottomrule
\end{tabular}%
}
\vspace{8pt}
\caption{We involve KIT~\cite{Plappert2016kit}dataset and evaluate the methods on the text-driven motion generation task. Please refer to Tab. 3 for more details on metrics and notations.
% ($cf.$\cref{tab:tm:comp:humanml3d}for metrics details)
}
\label{tab:tm:comp:kit}
\end{table}
     
\subsection{Ablation on Motion Tokenizer.} 
\label{sec:vqvae}
We ablate the motion tokenizer $\mathcal{V}$ of our MotionGPT models, studying the size $K$ of motion codebooks. We also compare this VQ-VAE with other VAE models in previous works~\cite{vposer_SMPL-X:2019, petrovich21actor, chen2023mld}, as shown in \cref{tab:mr:ablation}. This comparison demonstrates the improvement of VQ-VAE on motion reconstruction. With this ablation studies on the codebook size $K$, we thus select $K=512$ for most experiments.



\begin{table}[t]
\centering
\resizebox{0.7\columnwidth}{!}{%
\begin{tabular}{@{}lccccc@{}}
\toprule
\multirow{2}{*}{Method}&\multicolumn{5}{c}{Reconstruction}
\\\cmidrule(lr){2-6}
&MPJPE$\downarrow$&PAMPJPE$\downarrow$&ACCL$\downarrow$&FID$\downarrow$&DIV$\rightarrow$
\\ \midrule
Real  &  \multicolumn{1}{c}{-} & \multicolumn{1}{c}{-}& \multicolumn{1}{c}{-} &  $0.002\hidden{^{\pm.000}}$&	$9.503\hidden{^{\pm.065}}$     \\ \midrule
VPoser-t~\cite{vposer_SMPL-X:2019}  & $75.6$ & $48.6$& $9.3$ &     $1.430\hidden{^{\pm.002}}$ & $8.336\hidden{^{\pm.099}}$   \\
ACTOR~\cite{petrovich21actor}     & $65.3$ & $41.0$& $\boldsymbol{7.0}$ &    $0.341\hidden{^{\pm.002}}$ & $\boldsymbol{9.569}\hidden{^{\pm.099}}$ \\
MLD-1~\cite{chen2023mld}  & $\boldsymbol{54.4}$ & $41.6$ & $8.3$ & $0.247\hidden{^{\pm.001}}$ & $9.630\hidden{^{\pm.081}}$ \\
 \midrule
% 7 256
MotionGPT (Ours) & 
$55.8$ &$\boldsymbol{40.1}$  & $7.5$	&$\boldsymbol{0.067}\hidden{^{\pm.001}}$&$9.675\hidden{^{\pm.064}}$	
 % $57.5$ & $40.2$ & $7.3$ & $0.092\hidden{^{\pm.001}}$ & $9.724\hidden{^{\pm.068}}$ 
\\ \bottomrule
\toprule
$K=256$ & $76.4$ & $51.3$ & $10.0$ & $0.187\hidden{^{\pm.002}}$ & $\boldsymbol{9.496}\hidden{^{\pm.070}}$ \\ 
$K=512$ & $\boldsymbol{55.8}$ &$\boldsymbol{40.1}$  & $\boldsymbol{7.5}$	&$\boldsymbol{0.067}\hidden{^{\pm.001}}$&$9.675\hidden{^{\pm.064}}$ \\
$K=1024$ & $60.3$ & $44.0$ & $8.6$ & $0.086\hidden{^{\pm.001}}$ & $9.677\hidden{^{\pm.078}}$\\
$K=2048$ & $78.9$ & $51.4$ & $10.5$ & $0.141\hidden{^{\pm.001}}$ & $9.597\hidden{^{\pm.056}}$
% \\\midrule
% $K=1024,\text{Depth}=3, \text{Width}=512$ & $60.365$ & $44.087$ & $8.677$ & $0.086^{\pm.001}$ & $9.677^{\pm.078}$\\
% $K=1024,\text{Depth}=4, \text{Width}=512$ & $72.321$ & $48.797$ & $9.830$ & $0.062^{\pm.001}$ & $9.530^{\pm.077}$   \\
% $K=1024,\text{Depth}=5, \text{Width}=512$ & $77.370$ & $50.605$ & $10.263$ & $0.047^{\pm.001}$ & $9.544^{\pm.071}$
\\ \bottomrule
% \toprule
% \multirow{2}{*}{Method}&\multicolumn{3}{c}{Text-to-Motion} & \multicolumn{2}{c}{Motion-to-Text} \\
% \cmidrule(lr){2-4}\cmidrule(lr){5-6}
% &R@3$\uparrow$&FID$\downarrow$&DIV$\rightarrow$&Cider$\uparrow$&BertScore$\uparrow$
% \\ 
% \midrule
% MLD-1~\cite{chen2023mld} 
% \\\midrule
% CodeBook Size 512\\
% CodeBook Size 1024 \\
% Down Sample 2x\\
% Width 512\\
% Width 1025\\
% \midrule
% MotionGPT (Ours) \\
% \bottomrule
\end{tabular}%
}
\vspace{10pt}
\caption{Evaluation of our motion tokenizer on the motion part of HumanML3D~\cite{Guo_2022_CVPR_humanml3d} dataset. We follow MLD~\cite{chen2023mld} to evaluate our VQ-VAE model $\mathcal{V}$: MPJPE and PAMPJPE are measured in millimeter. ACCL indicates acceleration error. We evaluate FID and Diversity the same as Tab. 3. The baselines of VPoser-t~\cite{vposer_SMPL-X:2019} and ACTOR~\cite{petrovich21actor} are borrowed from MLD. $K$ indicates the codebook size, and $K=512$ shows the best performance of motion reconstruction.}
% \vspace{-10pt}
\label{tab:mr:ablation}
\end{table}

\newpage
% \subsection{Effectiveness of Pre-training.} 
\label{sec:pretrian}
\subsection{Effectiveness of Training Scheme} 
\textbf{Motion-Language Pre-training} vs \textbf{Instructions Tuning}. We have provided the illustration of our training scheme in Fig. 3 and the evaluation in Tab. 6. We further ablate this training scheme on the base MotionGPT model, by evaluating the motion-language pre-training (the second step) and instruction tuning (the third step). As shown in \cref{tab:abl:pretrain}, we train these models with the same 600K iterations. Compared to other training combinations, the full-stage MotionGPT achieves higher performance on most motion tasks.

\begin{table}[h]
\resizebox{\columnwidth}{!}{%
\begin{tabular}{@{}lcccccccccccc@{}}
\toprule
\multirow{2}{*}{Size} & 
\multirow{2}{*}{Pre-training} & 
\multirow{2}{*}{Instruction Tuning} & 
\multicolumn{3}{c}{Text-to-Motion}& 
\multicolumn{3}{c}{Motion-to-Text} & 
\multicolumn{2}{c}{Motion Prediction}&
\multicolumn{2}{c}{Motion In-between}
\\ \cmidrule(lr){4-6} \cmidrule(lr){7-9} 
\cmidrule(lr){10-11} 
\cmidrule(lr){12-13}
&&& R TOP3 $\uparrow$ & FID $\downarrow$ & DIV $\rightarrow$ & MMDist$\downarrow$ & Bleu@4$\uparrow$ & Cider$\uparrow$ 
&   FID $\downarrow$ & DIV $\rightarrow$  
& FID $\downarrow$ & DIV $\rightarrow$ 
\\ \midrule
Real & - & - &
$0.797\hidden{^{\pm.002}}$ &
$0.002\hidden{^{\pm.000}}$ &
$9.503\hidden{^{\pm.065}}$ &
% - &
% $0.828$ &
${2.901}$
& -& -
& $0.002$ & $9.503$ 
& $0.002$ & $9.503$ 
\\ \midrule
Base & \CheckmarkBold & \XSolidBrush&
$\boldsymbol{0.722}\hidden{^{\pm.002}}$	& $0.365\hidden{^{\pm.008}}$	& 
$9.407\hidden{^{\pm.042}}$ &
$\boldsymbol{2.821}$	& $\boldsymbol{12.47}$ & $\boldsymbol{29.2}$ 
& - & - 
& - & - 
\\
Base & \CheckmarkBold &  \CheckmarkBold &
$0.700$	& $\boldsymbol{0.160}$	& $\boldsymbol{9.411}$ &
$3.019$	& $11.42$ & $28.2$ 
% &0.668&9.172&0.157&9.389
& ${0.905}$ &${8.972}$
&$\boldsymbol{0.214}$&$\boldsymbol{9.560}$
\\
Base & \XSolidBrush  & \CheckmarkBold & 0.607 & 0.324 & 9.563 & 3.374 & $10.92$& $27.7$&  $1.643$ & $8.829$ & $0.323$  & $9.628$
\\ 
\bottomrule
\end{tabular}%
}
\vspace{10pt}
\caption{Evaluation of the training scheme on the base MotionGPT models. We evaluate the results with the proposed evaluation protocols in \cref{sec:appendix:data: details}. Please refer to Tab. 2 for metrics and the details.}
\label{tab:abl:pretrain}
\end{table}

\label{sec:finetuning}
\textbf{Instructions Tuning} vs \textbf{Task-Specific Tuning}. While our unified instruction-tuned MotionGPT model has demonstrated competitive performance across various motion-related tasks, further fine-tuning can always enhance its performance on specific tasks. Therefore, we focus on the text-to-motion task and motion in-between task as illustrative examples to showcase the performance of the model before and after fine-tuning. By comparing the results in \cref{tab:tm:abl:finetune}, we can assess the effectiveness of fine-tuning in improving task-specific performance.


\begin{table}[h]
\centering
\resizebox{0.8\columnwidth}{!}{%
\begin{tabular}{@{}cccccccccc@{}}
\toprule
% \multirow{2}{*}{Pretrained} & 
\multirow{2}{*}{Insturct tuned} & 
\multirow{2}{*}{Fine tuned} & & 
% \multirow{2}{*}{FT} & 
\multicolumn{3}{c}{Text-to-Motion}& & 
% \multicolumn{3}{c}{Motion-to-Text} 
% \multicolumn{2}{c}{Motion Prediction}&
\multicolumn{3}{c}{Motion In-between}
\\ \cmidrule(lr){4-6}  \cmidrule(lr){8-10} 
&& 
& R TOP1$\uparrow$ 
& FID$\downarrow$ 
& DIV$\rightarrow$ 
&&$\text{FID}\downarrow $  
& DIV$\uparrow$            
& ADE$\downarrow$
% & R TOP3 & Cider & BertScore  
% &  APD & ADE  &  APD & ADE 
\\ \midrule
 % &  & - & - & - \\
\CheckmarkBold & \XSolidBrush & &
${0.435}\hidden{^{\pm.003}}$ & $\boldsymbol{0.160}\hidden{^{\pm.008}}$ & $9.411\hidden{^{\pm.081}}$ 
&&
${0.214}$&$\boldsymbol{9.560}$&${3.762}$
\\
\CheckmarkBold & \CheckmarkBold  & &
$\boldsymbol{0.492}\hidden{^{\pm.003}}$ &
${0.232}\hidden{^{\pm.008}}$ & 
$\boldsymbol{9.528}\hidden{^{\pm.071}}$ &
&$\boldsymbol{0.209}$&$9.378$&$\boldsymbol{3.281}$
\\ \bottomrule
\end{tabular}%
}
\vspace{10pt}
\caption{Evaluation of new task tuning of different size models on HumanML3D~\cite{Guo_2022_CVPR_humanml3d} dataset. 
}
\label{tab:tm:abl:finetune}
\end{table}






\subsection{Difference of T2M-GPT}
\label{sec:appendix:diff}
We introduce the difference between T2M-GPT~\cite{zhang2023generating} to show our unified framework. T2M-GPT investigates a generative framework based on VQ-VAE and Transformer for motion generation only. 
%
They incorporate language information by leveraging CLIP \cite{radford2021clip} to extract text embedding as motion generation conditions, which is similar to most previous work, such as MDM~\cite{mdm2022human}, MLD~\cite{chen2023mld}, and MotionDiffuse~\cite{zhang2022motiondiffuse}.
%
% as the  text condition of their GPT model and train the GPT model from scrach. 
%
However, our MotionGPTs are based on the pre-trained language model so it naturally leverages the strong language generation and zero-shot transfer abilities of pre-trained language models. Benefiting from the motion-language vocabulary, MotionGPT thus generates both human language and human motion in a unified model.

\subsection{Implementation details of Motion Completion}
\label{sec:appendix:reimplement}
Please note that MDM\cite{mdm2022human} accomplish motion in-between task in their paper through masked motion “in-painting” which fix the first and last 25\% of the motion, leaving the model to generate the remaining 50\% in the middle. To achieve the motion prediction task with MDM, we fix the first 20\% of the motion and then generate the remaining. All our results are computed by utilizing their provided pre-trained model. To compare with MDM in Tab. 5 on both motion in-between and motion prediction tasks, we evaluate our MotionGPT with the same setting during the inference.


\section{Inference Time}
\label{sec:inferencetime}
%
We provide a detailed study on inference time with our different model sizes below. Due to our auto-regressive model for motion generation, we use Frames Per Second (FPS) to evaluate our time costs. All the time costs are evaluated on 8 Tesla V100 using one batch size. \cref{tab:tm:abl:scheduler} shows that any size of our MotionGPTs can support real-time human animations and come up to hundreds of FPS.

% avg gt frames 137.77
\begin{table}[h]
\centering
% \vspace{-6pt}
\resizebox{0.45\columnwidth}{!}
{%
\begin{tabular}{@{}llcc@{}}
\toprule
Models & Backbone&Parameters
% & FLOPs (G) $\downarrow$ 
& FPS $\uparrow$
\\ \midrule
% MDM \\ 22.8 & & 5.57
% MLD & 26.9 & 29.86 & 634.88
% \\ \midrule
MotionGPT & Small & 60 M 
% & 14.84 
& 421.31
\\
MotionGPT &Base & 220 M
% &  59.04 
& 222.69
\\
MotionGPT &Large & 770 M
% & 208.89 
& 119.75
\\
\bottomrule
\end{tabular}%
}
\vspace{5pt}
\caption{Evaluation of inference time costs on text-driven motion generation. We evaluate the Frames Per Second (FPS) by averaging our generated frames for each second. We show the time costs on different model sizes. Under the same 1 Tesla V100, the smaller model size gets the faster FPS. All models can support real-time motion animation applications.}
\vspace{-5pt}
\label{tab:tm:abl:scheduler}
% \vspace{-10pt}
\end{table}


\section{Statistics on Motion Vocabulary}
 % Utilization Rating
\label{sec:appendix:visual}
We visualize the usage of each "word" in our motion vocabulary $V_m$ item generated by our motion tokenizer $\mathcal{V}$. We sample all motions from the whole test set of HumanML3D dataset~\cite{Guo_2022_CVPR_humanml3d} and count each "word". In \cref{fig:codebook}, it shows the utilization of our motion codebook, which seems to be a concise but informative motion vocabulary.


\begin{figure}[!h]
\centering
\begin{minipage}[t]{0.32\textwidth}
\centering
\includegraphics[width=5cm]{codeBar_256.pdf}
% \caption{256}
\end{minipage}
\begin{minipage}[t]{0.32\textwidth}
\centering
\includegraphics[width=5cm]{codeBar_512.pdf}
% \caption{512}
\end{minipage}
\begin{minipage}[t]{0.32\textwidth}
\centering
\includegraphics[width=5cm]{codeBar_1024.pdf}
% \caption{1024}
\end{minipage}
\caption{The statistics of each "word" in different sizes of motion vocabulary $V_m$. From left to right, the vocabulary size is $K=256, 512,1 024$. ($cf$. \cref{tab:mr:ablation}, $K=512$ for the best motion quality.) }
\label{fig:codebook}
\end{figure}

% \begin{figure}[htb]
%     \centering
%     % \vspace{-10pt}
%     \includegraphics[width=1.0\linewidth]{rebuttal/TSNE.pdf}
%     % \vspace{-18pt}
%     \caption{Visualization of the t-SNE results on evolved latent codes $\hat{z}_t$ during the reverse diffusion process (inference) on action-to-motion task.
%     %
%     $t$ is the diffusion step but ordered in the forward diffusion trajectory.
%     %
%     $\hat{z}_{t=49}$ is the initial random noise. $\hat{z}_{t=0}$ is our prediction.
%     %
%     We sample 30 motions for each action label.}
%     \label{fig:tsne}
%     % \vspace{-10pt}
% \end{figure}

\section{Evaluation of Hyperparameters}
\label{sec:appendix:hyperparameters}
We conduct experiments to investigate the impact of different sampling strategies on the generation results. Specifically, we compare the use of greedy search, which selects the most probable token at each step, with sampling from the probability distribution and adopting beam search, which is evaluated in previous language models~\cite{raffel2020t5}. Beam search expands the search space for improved sequence probability matching. The results in \cref{tab:appendix:guidance} demonstrate that while avoiding sampling and using beam search can slightly improve generation quality, they also significantly reduce the diversity of generated motions from the same text description.

\label{appendix:guidance}

\begin{table}[H]
\centering
\resizebox{\columnwidth}{!}{%
\begin{tabular}{@{}lcccccccc@{}}
\toprule
 Method & \multirow{2}{*}{Sample} & \multirow{2}{*}{$\text{\#beams}$}  & \multicolumn{1}{c}{R Precision} & \multicolumn{1}{c}{\multirow{2}{*}{FID$\downarrow$}} & \multirow{2}{*}{MM Dist$\downarrow$} & \multirow{2}{*}{Diversity$\rightarrow$} & \multirow{2}{*}{MModality$\uparrow$} \\\ &&&  \multicolumn{1}{c}{Top 3$\uparrow$}       & \multicolumn{1}{c}{}                     &                          &                            &     \\ \toprule
Real & - &- &
  $0.797^{\pm.002}$ &
  $0.002^{\pm.000}$ &
  $2.974^{\pm.008}$ &
  $9.503^{\pm.065}$ &
  \multicolumn{1}{c}{-}
  \\ \midrule
\multirow{4}{*}{MotionGPT}  && - & $0.780^{\pm.002}$ & $0.224^{\pm.009}$ & $3.076^{\pm.009}$ & $9.492^{\pm.056}$ & \multicolumn{1}{c}{-}
\\
&&2 & $0.780^{\pm.002}$ & $0.199^{\pm.008}$ & $3.083^{\pm.007}$ & $9.512^{\pm.063}$ & \multicolumn{1}{c}{-}
\\
&&3 & $0.781^{\pm.002}$ & $0.179^{\pm.008}$ & $3.099^{\pm.009}$ & $9.516^{\pm.064}$ &  \multicolumn{1}{c}{-}
\\
&&4 & $0.782^{\pm.002}$ &	$0.160^{\pm.007}$ &	$3.092^{\pm.010}$ &	$9.536^{\pm.060}$	&  \multicolumn{1}{c}{-}

\\ \midrule
\multirow{4}{*}{MotionGPT} &\checkmark & - & ${0.778}^{\pm.002}$ & $0.232^{\pm.008}$ & ${3.096}^{\pm.008}$ & ${9.528}^{\pm.071}$ & ${2.008}^{\pm.084}$
\\
&\checkmark & 2 & $0.780^{\pm.002}$ & $0.194^{\pm.008}$ & $3.091^{\pm.010}$ & $9.508^{\pm.063}$ & {$1.140^{\pm.064}$}
\\
&\checkmark & 3 & $0.780^{\pm.002}$ & $0.190^{\pm.008}$ & $3.089^{\pm.011}$ & $9.529^{\pm.061}$ & $0.929^{\pm.055}$
\\
&\checkmark & 4 & $0.780^{\pm.002}$ & $0.182^{\pm.008}$ & $3.093^{\pm.008}$ & $9.537^{\pm.059}$ & $0.803^{\pm.044}$
\\
\bottomrule
\end{tabular}%
}
% \vspace{-6pt}
\vspace{5pt}
\caption{Evaluations on hyperparameters for MotionGPT generations. We study the influence of two hyperparameters: \textit{sample} stands for sampling from distribution; $\textit{\#beams}$ means the number of beams for beam search, where empty means no beam search.}
\label{tab:appendix:guidance}
% \vspace{-6pt}
\end{table}



\section{User Study}
\label{sec:appendix:userstudy}
For the comparisons of text-to-motion task, we use the force-choice paradigm to ask ``Which of the two motions is more realistic?'' and ``which of the two motions corresponds better to the text prompt?''. The provided motions are generated from 30 text descriptions from the test set of HumanML3D~\cite{Guo_2022_CVPR_humanml3d} dataset. For the comparisons of motion-to-text task, we ask 15 users to choose the motion descriptions from GT, TM2T~\cite{chuan2022tm2t}, and our MotionGPT. The motions are from the test set of HumanML3D~\cite{Guo_2022_CVPR_humanml3d} dataset.
%
As shown in \cref{fig:appendix:user}, in both two tasks, our MotionGPT was preferred over the other state-of-the-art methods and even competitive with the ground truth.
%

% \begin{figure}[!h]
% \centering
% \begin{minipage}[t]{0.65\textwidth}
% \centering
% \includegraphics[width=9cm]{user_study_1.pdf}
% % \caption{256}
% \end{minipage}
% \begin{minipage}[t]{0.34\textwidth}
% \centering
% \includegraphics[width=4cm]{user_study_2.pdf}
% % \caption{512}
% \end{minipage}
% \caption{Visualizations of the usage of each "word" in different sizes of motion vocabulary $V_m$ item. From left to right the vocabulary size is $K=256,512,1024$.}
% \label{fig:appendix:user}
% \end{figure}
\begin{figure}[!h]
\centering
\includegraphics[width=\textwidth]{user_study.pdf}
\vspace{-10pt}
\caption{User Study. We investigate our motion quality and the alignment with test descriptions. The left part is the user study for text-to-motion. The right part is for motion captioning.}
\label{fig:appendix:user}
\end{figure}


\section{Evaluation Protocols on the Uniform Motion-Language Generation.}
\label{sec:appendix:data: details}
We propose a protocol to evaluate our unified MotionGPT on multiple motion-language generation tasks. Upon previous datasets~\cite{Guo_2022_CVPR_humanml3d,Plappert2016kit,AMASS_ICCV2019}, we build an instruction motion-language dataset, which is composed of 14 core tasks (\cref{fig:appendix:dataset}) for now. As shown in \cref{tab:appendix:prompts}, each core task has dozens of instruction prompts (\cref{tab:appendix:prompts}). We will release the pre-processed dataset.

\begin{figure}[h]
    \centering
    \includegraphics[width=\linewidth]{instructdata.png}
    \caption{Protocols for multiple motion-language tasks. For each task, we follow \cref{tab:appendix:prompts} to process the previous datasets~\cite{AMASS_ICCV2019, Guo_2022_CVPR_humanml3d} into the instruction-based data.}
    \label{fig:appendix:dataset}
\end{figure}

\begin{table}[H]
\centering
\resizebox{\columnwidth}{!}{%
\begin{tabular}{@{}lll@{}}
\toprule
Task & Input & Output
\\ \midrule
\multirow{3}{*}{Text-to-Motion}& Give me a motion that corresponds to [caption]. &\multirow{3}{*}{[motion] }
\\ & Demonstrate a sequence of movements that depict [caption]. &
\\ & I need a human motion that conveys [caption]. Can you generate it for me?&
\\ \midrule \multirow{2}{*}{Text-to-Motion w/ length}& Give me a motion that lasts for approximately [frames] frames. The caption is: [caption]. &\multirow{2}{*}{[motion] }
\\ & Please create a motion that lasts [seconds] seconds and illustrates [caption]. &
\\ \midrule \multirow{2}{*}{Length-to-Motion }& Show me a motion that lasts for no more than [frames] frames. &\multirow{2}{*}{[motion] }
\\ & Create a motion that has a duration of [seconds] seconds. &
\\ \midrule \multirow{2}{*}{Radnom Motion }& Give me motions as you like. &\multirow{2}{*}{[motion] }
\\ & Produce actions that are not prescribed. &
\\ \midrule \multirow{2}{*}{Motion-to-Text}& Give me a summary of the motion being displayed in [motion] using words. &\multirow{2}{*}{[caption] }
\\ & Describe the motion illustrated in [motion] in natural language. & 
\\ \midrule \multirow{2}{*}{Motion-to-Text w/ length}& Describe the movement portrayed in [motion]that lasts [frames] frames. &\multirow{2}{*}{[caption] }
\\ & What is happening in [motion] for a length of [seconds] seconds? &
\\ \midrule \multirow{2}{*}{Motion-to-Length}& What is the duration of [motion]'s gestures in frames? & There are [frames] frames in the motion.
\\ & What is the total duration of  [motion]'s body movements in seconds? & The motion lasts for [seconds] seconds.
\\ \midrule \multirow{2}{*}{Caption-to-Length}&How many frames are expected for the motion that matches [caption]?& The duration is estimated to be around [frames] frames.
\\ & Given [caption], provide the anticipated second duration for the corresponding motion. & The motion has a length of [seconds] seconds.
\\ \midrule \multirow{2}{*}{Length-to-Caption}&What are some possible physical gestures that could be made in [frames] frames?&\multirow{2}{*}{[caption] }
\\ & What motion could be performed in [seconds] seconds? &
\\ \midrule \multirow{2}{*}{Random Caption}&Depict a motion as like you have seen it.&\multirow{2}{*}{[caption] }
\\ & Describe the motion of someone randomly.&
\\ \bottomrule
\end{tabular}%
}
\vspace{10pt}
\caption{Some examples of prompt templates in our uniform evaluation protocols.}
\label{tab:appendix:prompts}
\end{table}


% \section{Metric Definitions}
\label{appedix:metrics:details}
% AITS 详细计算流程


% motion quality, generation diversity, condition matching, and time costs.
\textbf{Metric Definitions:} We provide more details of evaluation metrics as follows. Our evaluation metrics can roughly divide to five classes including text-motion matching, generation diversity, linguistic quality, motion quality, and time cost. For the first two classes, \cite{chen2023mld} has already claims clearly and for the linguistic metrics including BLUE~\cite{papineni2002bleu}, Rouge~\cite{lin2004rouge}, Cider~\cite{vedantam2015cider}, and BertScore~\cite{zhang2019bertscore}, you can refer to their own papers for details. Here we focus on the explaination of the rest metrics.

% \textbf{NLP Metrics}. xxxxx 

% \cx{xx R-Percision xx  Multimod. xxx MM Dist xxKID xx}
\textbf{Motion Quality}. 
FID, MPJPE, PAMPJPE~\cite{gower1975generalized}, ACCL have been clearly explained in \cite{chen2023mld}. Thus here we focus on the  Average Displacement Error (ADE) and Final Displacement Error (FDE) refaccuracy of the predicted motion. Following previous motion prediction work\cite{yuan2020dlow, zhang2021we, ma2022multi}, ADE is defined as average L2 distance between the ground truth and predicted motion of the whole sequence and FDE is the L2 distance between the ground truth and predicted motion in the last frame.

\textbf{Time Costs}. To evaluate the computing efficiency of our models, especially the inference efficiency, we calculate average Frames Per Second (FPS) when generating motions.  In our case, we calculate FPS on the test set of HumanML3D~\cite{Guo_2022_CVPR_humanml3d}, set the batch size to one, and ignore the time cost for model and dataset loading parts.


\section{Details on MotionGPT Models}
\label{sec:appendix:model:details}


\subsection{Implementation Details}
Besides the MotionGPt with 220M parameters, we implement a smaller model that reduces the model dimension with $d_{\text{model}} = 512$, $d_\text{ff} = 2048$ with only 6 layers in encoder and decoder, as well as a larger model with 770 million parameters, which increases the model dimensions with $d_\text{model} = 1024$, $d_\text{ff} = 4096$, $d_\text{kv} = 64$, 24 layers for each transformer. Except for the training iterations during the instruction tuning stage, the other settings are the same. Please refer to \cref{tab:appendix:details} for more details.


\begin{table}[H]
\centering
\resizebox{0.65\columnwidth}{!}{%
\begin{tabular}{@{}lccc@{}}
\toprule
MotionGPT &  Small & Base & Large
\\ \toprule
Backbone & Flan-T5-Small & Flan-T5-Base & Flan-T5-Large \\
Training Batch Size & 64 &16 &4 
\\
Model Size& 60M&220M&770M
\\
Pre-training - Iterations & 300K& 300K& 300K
\\
Pre-training - Learning Rate & 2e-4& 2e-4& 2e-4
\\
Instruction Tuning - Iterations & 200K& 300K& 400K
\\
Instruction Tuning - Learning Rate & 1e-4& 1e-4& 1e-4
\\
\midrule
Motion Vocabulary Number $V_m$ & 512 & 512 & 512 \\
Motion Codebook Dimension & 512 & 512 & 512\\
% Motion Tokenizer $\mathcal{V}$ Depth & 3 & 3 & 3 
% \\
% Motion Tokenizer - Training Batch Size & 256 & 256 & 256
% \\
% Motion Tokenizer - Training Iterations & 150K& 150K& 150K
% \\
% Motion Tokenizer - Learning Rate & 1e-4& 1e-4& 1e-4
% \\
\bottomrule
\end{tabular}%
}
\vspace{10pt}
\caption{Hyperparameters for different MotionGPTs. We train these models on 64 Tesla V100 GPUs. The smaller model size lowers the computational requirements and thus provides faster inference ($cf.$ \cref{sec:inferencetime}). According to Tab. 6, the base MotionGPT model is the best one for overall tasks. However, we believe this could be caused by the small amount of current motion datasets. The large model could achieve the best performance when the amount of data comes up to millions or even billions. }
\label{tab:appendix:details}
\end{table}


% % Methods details
% \subsection{Details Information on Variational Autoencoder Models}
% \label{sec:appendix:loss}
% %We train Variational Autoencoder Models $\mathcal{V}$ on .
% \jb{We take HumanML3D~\cite{Guo_2022_CVPR_humanml3d} and its motion representation (\cref{sec:appendix:motionRepre}) as an example here to explain our loss details of Variational Autoencoder Models $\mathcal{V}$.}
% %
% % The models is optimized with two loss terms.
% %
% The motion $x^{1:L}$ includes joint features and is supervised with data term by mean squared error:
% \begin{align}
% % (x^{1:L}, \mathcal{D}(\mathcal{E}(x^{1:L})))
% \mathcal{L}_{data} = \left\| x^{1:L} - \mathcal{D}(\mathcal{E}(x^{1:L})) \right\|^2.
% %
% \end{align}
% To regularize latent space as a standard variational autoencoder~\cite{kingma2013auto}, we employ a Kullback-Leibler term between $q(z|x^{1:L}) = \mathcal{N}\left(z ; \mathcal{E}_\mu, \mathcal{E}_{\sigma^2}\right)$ and a standard Gaussion distribution $\mathcal{N}\left(z;0,1\right)$.
% %
% The full training loss of the VAE model $\mathcal{V}$ follows:
% \begin{align}
% \mathcal{L}_{\mathcal{V}} = \mathcal{L}_{data}(x^{1:L}, \mathcal{D}(\mathcal{E}(x^{1:L}))) + \lambda_{reg}\mathcal{L}_{reg}(x^{1:L};\mathcal{E},\mathcal{D}),
% \end{align}
% where $\lambda_{reg}$ is a low weight to control the regularization. 
% %
% The KIT~\cite{Plappert2016kit}, HumanAct12~\cite{guo2020action2motion} and UESTC~\cite{ji2018large} dataset processed by \cite{petrovich22temos,petrovich21actor} also supports SMPL-based~\cite{SMPL2015} motion representation.
% %
% Here we list the loss terms for this representation.
% %
% The data term formulates as followed:
% \begin{align}
%     \label{eq:vae:humanml}
%     \mathcal{L}_{data} = \sum_{i=1}^L \left\|{r}^i  -\hat{{r}}^i \right\|_2 + \sum_{i=1}^L \left\|{\theta}^i  -\hat{{\theta}}^i \right\|_2 +  \left\|{\beta}  -\hat{{\beta}} \right\|_2.
% \end{align}
% Here the motion is $x^{1:L} = \{{r}^i, {\theta}^i, {\beta}\}^{L}_{i=1}$, which includes global translation ${r}^i$, pose parameter ${\theta}^i$ and shape parameter ${\beta}$ of the $i$-th frame. 
% %
% To enhance the full-body supervision, the reconstruction term on the SMPL vertices follows:
% \begin{align}
%     \mathcal{L}_{\mathrm{mesh}} = \sum_{i=1}^L \left\| {V}_i - M(\hat{{r}}^i, \hat{{\theta}}^i,\hat{{\beta}}^i) \right\|^2,
% \end{align}
% where the body reconstruction function $M(\cdot)$ is from the differentiable SMPL layer, while the vertices ${V}_i$ are calculated with the ground truth motion parameters using the same layer.
% %
% The reconstruction loss builds global supervision on almost all predicted parameters $ \{{r}_t, {\theta}_t, {\beta}\} $ and shows a reliable supervision~\cite{petrovich21actor} for motion generation.
% %
% The full objective on SMPL-based motion representation reads:
% \begin{align}
%     \label{eq:vae:smpl}
%     \mathcal{L}_{\mathcal{V}} = \mathcal{L}_{data} + \lambda_{mesh}\mathcal{L}_{mesh} + \lambda_{reg}\mathcal{L}_{reg}.
% \end{align}
% % 总体解释，每个loss的目的和作用
% where $\lambda_{mesh}$ is the weight to enhance the supervision on the full-body vertices. Besides, the regularization term is the same as the Kullback-Leibler term in \cref{eq:vae:humanml}.
% %
% In practice, the shape parameters, as part of global motion features, increase the complexity of motion generation and influence joint positions. 
% % Thus, the shape error causes the foot sliding easily. 
% % \cref{sec:appendix:motionRepre}, we 
% We finally utilize the objective of ~\cref{eq:vae:humanml} to train our text-based models and \cref{eq:vae:smpl} to train action-based models in comparisons and evaluations.


% \subsection{Network Architectures}
% The details of network architecture are shown as \cref{fig:appendix:net}, our MLD comprises three main components, motion encoder $\mathcal{E}$, motion decoder $\mathcal{D}$ and latent denoiser $\epsilon_\theta$. Please refer to the following figure and \cref{tab:appendix:details} for more details.
% % 2.1 network structure
% % 2.3 implement details

% \begin{figure}[h]
%     \centering
%     \includegraphics[width=0.85\linewidth]{network.pdf}
%     \caption{Network architecture of our conditional MLD. We explain each component in its right bottom part and the loss terms in \cref{sec:appendix:loss}.  }
%     \label{fig:appendix:net}
% \end{figure}


